\documentclass[tikz,border=22pt]{standalone}
\usepackage[UTF8]{ctex}
\usepackage{tikz}
\usetikzlibrary{arrows.meta,positioning,shapes.geometric,shapes.misc,fit,calc,shadows,backgrounds}

\begin{document}
\sffamily

% =================================================================
% 调色板定义 (Obsidian Excalidraw 原生高质感深色模式配色)
% =================================================================
\definecolor{canvasBg}{RGB}{24, 24, 24}
\definecolor{benchBorder}{RGB}{46, 180, 100}
\definecolor{benchBg}{RGB}{30, 42, 35}

\definecolor{agentBorder}{RGB}{235, 95, 90}
\definecolor{agentBg}{RGB}{48, 28, 30}
\definecolor{agentText}{RGB}{250, 220, 220}

\definecolor{nodeBorder}{RGB}{200, 205, 210}
\definecolor{nodeBg}{RGB}{32, 34, 38}
\definecolor{nodeText}{RGB}{240, 243, 246}

\definecolor{accentGreen}{RGB}{80, 210, 130}
\definecolor{lineColor}{RGB}{205, 215, 225}

\pagecolor{canvasBg}

\begin{tikzpicture}[
    scale=1,
    >=Stealth,
    % 基础节点样式
    rectNode/.style={
        rectangle, rounded corners=3pt, align=center,
        draw=nodeBorder, fill=nodeBg, text=nodeText,
        line width=1.3pt, inner sep=6pt, font=\small
    },
    agentStyle/.style={
        rectangle, rounded corners=8pt, align=center,
        draw=agentBorder, fill=agentBg, text=agentText,
        line width=1.8pt, inner sep=8pt, font=\bfseries\small
    },
    diamondNode/.style={
        diamond, aspect=1.65, align=center,
        draw=nodeBorder, fill=nodeBg, text=nodeText,
        line width=1.3pt, inner sep=3pt, font=\small
    },
    circleNode/.style={
        circle, align=center,
        draw=nodeBorder, fill=nodeBg, text=nodeText,
        line width=1.3pt, inner sep=3pt, font=\small
    },
    % 连线与标签样式
    flowLine/.style={
        ->, line width=1.2pt, draw=lineColor,
        >={Stealth[length=6.5pt, width=5pt]}
    },
    arcLine/.style={
        ->, line width=1.2pt, draw=lineColor,
        >={Stealth[length=6.5pt, width=5pt]}
    },
    textBadge/.style={
        fill=canvasBg, inner sep=2pt, font=\small, text=nodeText, rounded corners=2pt
    }
]

    % =============================================================
    % 节点坐标布局 (完全保真还原原图的物理与逻辑拓扑)
    % =============================================================

    % --- 1. Benchmark 顶层框内节点 ---
    \node[rectNode, minimum width=3.0cm, minimum height=1.35cm] at (-4.8, 3.8) (inst) {
        Instruction.md\\[2pt]
        \scriptsize (实验目标)
    };

    \node[agentStyle, minimum width=2.5cm, minimum height=1.55cm] at (-1.0, 3.8) (agent) {
        Agent\\[2pt]
        Loop
    };

    \node[diamondNode, minimum width=2.8cm, minimum height=1.55cm] at (2.8, 3.8) (eval) {
        内部评分机制
    };

    \node[circleNode, minimum size=1.35cm] at (5.8, 3.8) (score) {
        score
    };

    % Benchmark 容器框
    \begin{scope}[on background layer]
        \node[
            fit=(inst)(score),
            draw=benchBorder, dashed, line width=1.3pt, rounded corners=8pt,
            fill=benchBg!50,
            inner sep=14pt,
            label={[font=\bfseries\small, text=accentGreen, yshift=-2pt]above:\textbf{Benchmark}}
        ] (benchBox) {};
    \end{scope}

    % --- 2. 中间层：resource 与 evolve 目标 ---
    \node[rectNode, minimum width=3.0cm, minimum height=1.35cm] at (-1.0, 0.7) (resource) {
        \textbf{\textcolor{accentGreen}{resource}}
    };

    \node[rectNode, minimum width=2.8cm, minimum height=1.35cm] at (5.8, 1.3) (evolve) {
        \textcolor{accentGreen}{evolve} 目标
    };

    % --- 3. 诊断与改进层 ---
    \node[diamondNode, minimum width=3.2cm, minimum height=1.65cm] at (1.8, -1.4) (physician) {
        AI physician
    };

    \node[rectNode, minimum width=2.8cm, minimum height=1.35cm] at (5.8, -1.4) (action) {
        \textbf{\textcolor{accentGreen}{改进措施}}
    };

    % --- 4. 形式化验证与知识沉淀 (右翼) ---
    \node[diamondNode, minimum width=2.0cm, minimum height=1.3cm] at (9.4, -1.4) (lean) {
        lean
    };

    \node[rectNode, minimum width=2.5cm, minimum height=1.35cm] at (9.4, 0.7) (verifK) {
        verified\\[2pt]
        Knowledge
    };

    % --- 5. 技能工具沉淀 (左下翼) ---
    \node[rectNode, minimum width=3.6cm, minimum height=1.45cm] at (-2.0, -3.4) (skill) {
        skill / mcp /\\[2pt]
        knowledge card
    };

    % =============================================================
    % 连线系统 (严格保真原图逻辑流，自然优雅的贝塞尔样条曲线)
    % =============================================================

    % --- 1. Benchmark 内部主干流 ---
    \draw[flowLine] (inst) -- (agent);
    \draw[flowLine] (agent) -- (eval);
    \draw[flowLine] (eval) -- (score);

    % --- 2. score 驱动 evolve 目标 ---
    \draw[flowLine] (score) -- (evolve);

    % --- 3. resource 的三大汇聚输入 ---
    % (a) Agent Loop -> resource (垂直直连)
    \draw[flowLine] (agent.south) -- (resource.north);

    % (b) Instruction.md -> resource (弧线从底部向下并向右接入 resource 左侧)
    \draw[arcLine] (inst.south) to[out=270, in=180] (resource.west);

    % (c) 内部评分机制 -> resource (弧线从底部向下并向左接入 resource 右侧)
    \draw[arcLine] (eval.south) to[out=270, in=0] (resource.east);

    % --- 4. AI physician 的两大输入 ---
    % (a) resource -> AI physician (从 resource 底部弧线汇入 AI physician 左翼)
    \draw[arcLine] (resource.south) to[out=290, in=155] (physician.west);

    % (b) evolve 目标 -> AI physician (从 evolve 目标左侧弧线汇入 AI physician 顶部)
    \draw[arcLine] (evolve.west) to[out=180, in=75] (physician.north);

    % --- 5. AI physician -> 改进措施 ---
    \draw[flowLine] (physician) -- node[above, font=\small\bfseries, text=accentGreen] {总结} (action);

    % --- 6. 改进措施的两大分支 ---
    % 分支 A: 改进措施 -> lean (可形式化验证的) -> verified Knowledge -> Agent Loop
    \draw[flowLine] (action) -- node[above, font=\scriptsize, align=center, text=nodeText] {可形式化\\验证的} (lean);
    \draw[flowLine] (lean) -- (verifK);
    % verified Knowledge 大跨度优美弧线注入 Agent Loop
    \draw[arcLine] (verifK.north) to[out=95, in=20] ($(agent.north)+(0.3, 0)$);

    % 分支 B: 改进措施 -> skill / mcp / knowledge card (others) -> Agent Loop
    \draw[arcLine] (action.south) to[out=220, in=0] 
        node[textBadge, sloped, midway, above=1pt] {others} (skill.east);
    % skill / mcp / knowledge card 从外侧绕过 resource，在空旷处与 instruction 线优雅交叉，注入 Agent Loop 底部
    \draw[arcLine] (skill.north) to[out=135, in=270] (-3.4, 0.7) to[out=90, in=230] ($(agent.south)+(-0.35, 0)$);

\end{tikzpicture}
\end{document}